\begin{qfdefinition}[Macroeconomic indicators for macro trading: growth, inflation, employment, external accounts, fiscal, consumption, housing and business surveys, and demographics, with the FOMC's inflation goal and Gliner's demographic table read correctly]\label{def:macroeconomic-indicators}
\qffield{Statement} Gliner's chapter on data releases \citep[chap.~12, PDF pp.~401--428]{s7-gliner} groups the indicators a macro trader watches:
\begin{itemize}
\item \emph{Growth} (sec.\ ``Measuring Growth'', p.~402): GDP, nominal and real, quarterly with revisions.
\item \emph{Inflation} (sec.\ ``Inflation'', pp.~404--405): CPI and core CPI, PCE, the euro area's HICP, and the PPI, which measures ``the average change over time in the selling prices received by domestic producers''.
\item \emph{Employment} (sec.\ ``Employment and Population'', pp.~406--409): the unemployment rate, the labour force and participation rate (labour force as a share of the civilian noninstitutional population), and nonfarm payrolls, ``released on the first Friday of the month'', with average hourly earnings.
\item \emph{External accounts} (sec.\ ``Balance of Payments'', pp.~410--412): the current account (trade balance, net income and transfers), the trade balance, and the terms of trade, export prices relative to import prices.
\item \emph{Government} (sec.\ ``Government Indicators'', pp.~413--415): the budget balance and debt-to-GDP.
\item \emph{Consumption} (sec.\ ``Consumption Indicators'', pp.~416--417): U.S.\ consumption ``more than 70 percent of GDP in 2011'' and, with the United States at about 20 percent of world GDP, the U.S.\ consumer at ``a little more than 15 percent of global GDP'' (p.~416); retail sales, the savings rate and consumer confidence (Conference Board, University of Michigan).
\item \emph{Industry and services} (sec.\ ``Industry and Services Indicators'', pp.~418--420): purchasing managers' indices, with above 50 read as expansion and new orders and inventories among their components; industrial production; building permits, which lead housing starts; and capacity utilisation.
\item \emph{Demographics} (sec.\ ``Demographics'', pp.~421--426): working-age (15--64) and prime-saving (35--69) populations, with Table~12.2 (U.S.\ Census Bureau; heading and source on p.~424, data on p.~425) giving each country's peak year for both, and Trichet's remark that the current account ``is an important summary indicator'' (p.~424).
\end{itemize}

\qffield{Notation} \emph{Headline} and \emph{core} inflation (core excludes food and energy); PCE the personal consumption expenditures price index; PMI a diffusion index, $\%\text{better}+\tfrac12\%\text{same}$, scaled so that 50 means no change.

\qffield{Conditions} \emph{The Fed's inflation goal.} Gliner writes that ``the Fed targets core PCE'' (p.~405), and more broadly that central banks target core inflation. The FOMC's goal is 2 percent inflation ``as measured by the annual change in the price index for personal consumption expenditures'', the headline index; the statement was first adopted effective 24 January 2012 and was reaffirmed effective 27 January 2026 \citep[Statement on Longer-Run Goals and Monetary Policy Strategy, pp.~1--2]{fomc-longer-run-goals}. Core measures are watched as guides to the trend, not targeted. Gliner's further statements that PCE ``targets urban consumers'' and HICP ``rural consumers'' are wrong as descriptions: urban coverage defines the CPI-U, not PCE, and the HICP covers household consumption in general.

\emph{Payroll dates.} The Employment Situation is usually published on the first Friday but not always; the CES FAQ puts the release ``typically on the third Friday after the conclusion of the reference week'' \citep[Data Overview, Q8]{bls-ces-faq}, which usually, not always, falls on the first Friday of the next month; the BLS's own schedule is the record (the next release, as of October 2026, 6 November 2026) \citep[program page]{bls-ces-page}.

\qffield{Practice} \emph{Demographic peaks.} Gliner's two columns are easily mixed: Netherlands 1989, U.K.\ 2009 and U.S.\ 2007 are working-age (15--64) peaks, while Germany 2006, Japan 2016, China 2032, Korea 2023, Russia 2025, Spain 2022, India, Indonesia and the Philippines 2050 are prime-saving (35--69) peaks; Nigeria shows 2050 in both (Table~12.2, heading p.~424, data p.~425). These are projections published around 2013, and they have moved.

\emph{Against the UN's World Population Prospects 2024} \citep[file \nolinkurl{WPP2024_PopulationByAge5GroupSex_Medium}]{un-wpp-2024}: the 15--64 peaks are Germany 1997, Japan 1995, China 2015, Korea 2017, Russia 2010 and Spain 2009 (Gliner: 1986, 1992, 2011, 2015, 2011 and 2004), but the Netherlands 2026, the United Kingdom 2044, India 2048, Indonesia 2045 and the Philippines 2053, against his 1989, 2009, 2027, 2023 and 2037, and the United States and Nigeria are still rising at 2050 (the U.S.\ to 2100); the 35--69 peaks are Germany, Korea, Russia and Spain all in 2024, Japan 2010, China 2031, the Netherlands 2030, and after 2050 for the United Kingdom, the United States, India, Indonesia, Nigeria and the Philippines.

\emph{Why the years differ is not established here}. Gliner's table is labelled ``Source: U.S. Census Bureau'' (p.~424) with no vintage or release date; WPP 2024 is another publisher's 2024 revision. Publisher, vintage, base-period estimates and, for future years, projection assumptions all differ at once, and no controlled comparison was made (the same publisher across vintages, or both publishers in one vintage; the Census Bureau's International Database now requires an API key), so no cause is attributed. Two kinds of year should also be kept apart.

Peaks before about 2012 were historical estimates in both sources (WPP 2024 gives estimates to 2023 and projections after): there, a gap such as Germany's 15--64 peak in 1986 against 1997, or Japan's in 1992 against 1995, cannot come from projection assumptions and must reflect estimation, revision or territorial coverage (for example, how Germany before 1990 is defined). Peaks after about 2013 were projections in Gliner's table, and those after 2023 are still projections in WPP 2024; revised fertility, mortality and migration assumptions may contribute there, but that too is unverified. The point that projected peaks are revised with each release matters more than any single year.

\emph{Other claims.} Consumption ``more than 70 percent of GDP'' in the United States is Gliner's (chap.~8, p.~238; chap.~12, p.~416); U.S.\ consumer at about 15 percent of global GDP is also his (p.~416), the product of two approximate 2011 shares rather than a measured figure. PPI leading CPI by one to three months, that the ISM new-orders-minus-inventories spread is a reliable leading indicator, the quits rate as an indicator of labour-market tightness, wage growth driving stickier services inflation, median ages, home-price indices lagging the market, household formation as the structural driver of housing demand, the dependency ratio as a measure of fiscal and savings pressure, and a positive primary balance letting debt-to-GDP stabilise when growth exceeds the interest rate are conventional practitioner statements with no source here.

\qffield{Pitfalls} \emph{Capacity utilisation} ``over 100 percent'' (Gliner, p.~420) is not a reading one sees in the Federal Reserve's series, which measures output relative to sustainable capacity and has stayed well below 100 percent; read deviations from its own history. Debt-to-GDP above 90 percent as a threshold (Gliner, p.~414) is the contested Reinhart--Rogoff finding, not a rule.

\emph{Nonfarm payroll scope}. The BLS Handbook of Methods states that CES data ``exclude sole proprietors, the unincorporated self-employed, unpaid volunteer or family workers, farm workers, and domestic workers'', that ``Salaried officers of corporations are included'', and that government employment ``covers only civilian employees; military personnel are excluded'' \citep[CES concepts, page modified 28 February 2025]{bls-hom-ces}. The CES FAQ lists employees of non-profits and religious organisations as in scope, the largest group not covered by unemployment-insurance records, for which CES estimates separately \citep[Employment Concepts, Q4 and~Q6]{bls-ces-faq}. Nonprofit employees are therefore included; farm, private household, self-employed and military workers are excluded.

\qffield{See also} \hyperref[data:us-macro-release-calendar]{U.S.\ release calendar}; \hyperref[meth:macro-event-queries]{event queries}; \hyperref[meth:equity-macro-leading-indicators]{equity leading indicators}; \hyperref[def:tips-breakevens]{TIPS and breakevens}; \hyperref[def:central-bank-policy-framework]{central-bank policy}.

\qffield{Further reading} Nowcasting with release vintages \citep{wide-giannone-nowcasting}.

\end{qfdefinition}

\begin{qfmethod}[Leading indicators for equities in Gliner's macro toolkit: commodity prices, the CDX high-yield index, the ISM PMI, the Baltic Dry Index, consumer confidence and Australian-dollar implied volatility]\label{meth:equity-macro-leading-indicators}
\qffield{Statement} Gliner lists indicators that lead or confirm equity moves \citep[chap.~8, sec.\ ``Valuation Techniques for Equities'', PDF pp.~233--240]{s7-gliner}. The year-on-year change in commodity prices matters because a rise is inflationary and invites tightening (pp.~233--234). The five-year CDX high-yield index (CDX.NA.HY) rises with risk aversion and is inversely related to equities (pp.~234--235). For the ISM manufacturing PMI, when ``the ISM \dots\ goes below 50, the United States typically enters recession'' (p.~235). The Baltic Dry Index, a composite of Capesize, Panamax, Handysize and Supramax dry-bulk freight assessments, is watched because raw materials lead (p.~236).

His list also includes the ARMS index (pp.~237--238) and consumer confidence, since consumption is over 70 percent of U.S.\ GDP (p.~238). Three-month AUD/USD implied volatility is ``highly correlated'' with the VIX and inversely with the S\&P 500, and is most useful as a rate of change, though ``given the choice, using the VIX is a more effective factor'', and the two combined ``smooth out any error'' (pp.~239--240).

\qffield{Notation} CDX.NA.HY the Markit North American high-yield CDS index, quoted as a price or spread; ISM the U.S.\ Institute for Supply Management.

\qffield{Conditions} Gliner's charts to 2013. These are Gliner's chart readings (Figures~8.12--8.18), without lead-lag statistics or out-of-sample tests; the reading of ISM below 50 as recession has had false signals, since the index has fallen below 50 outside recessions.

\qffield{Practice} Treat each as a candidate nowcasting signal: build it from real-time vintages and test its lead against equity returns.

\qffield{Pitfalls} Neither the absence of a directly traded index future nor a slow supply response makes freight prices immune to speculation. Dry-bulk freight is traded through forward freight agreements settled on Baltic route assessments, and fleet supply changes with deliveries and scrapping.

\qffield{See also} \hyperref[def:macroeconomic-indicators]{macroeconomic indicators}; \hyperref[meth:equity-index-valuation]{equity-index valuation}; \hyperref[def:vix-variance-dividend-swaps]{VIX}.

\qffield{Further reading} Mixed-frequency nowcasting \citep{wide-banbura-nowcasting}.

\end{qfmethod}
